This activity covers user stories US-001, US-020, US-040, US-046 and US-048. It implements the schema designed in Activity~\ref{act:1.3}; this section describes only what the implementation adds to that design.

\subsection{Implementation of the schema}

The thirteen application tables of Figure~\ref{fig:er_diagram} are defined as Django models, the fourteenth being Django's own user table, and are created by seventeen migrations generated with \texttt{makemigrations} and kept under version control, so any installation reaches the same schema by running \texttt{migrate}. A check that reports no pending model changes runs before every story is closed. PostgreSQL 18 runs as a Docker Compose service with a health check, and the backend container applies the migrations before starting Gunicorn. SQLite is used only for fast local iteration; constraints and migrations are validated against the PostgreSQL service, because SQLite enforces some constraints differently.

The integrity rules of the design are declared on the models and created by the migrations as PostgreSQL constraints. Besides plain unique and check constraints, the schema uses a conditional unique constraint: a video can have only one statistics snapshot per collection run, while snapshots whose run was deleted are not restricted. A test asks the database itself whether every declared constraint exists, under its name and with its kind, so a constraint declared in Python but missing from the database would be detected.

Each of the twenty-one foreign keys has a deliberate deletion rule: sixteen cascade and five set the reference to null. The five nullable references keep methodological records alive when what they point to disappears: a run survives the deletion of its saved query or of the user who started it, a processing report survives its collection run, a snapshot survives its run and a video survives the purge of its channel. A test compares the rule of every foreign key with an expected table, so adding a key without deciding its rule fails the test suite.

\subsection{Shared resources and idempotent writes}

Videos, channels and comments are shared: a video found by two runs, or by two projects, is stored once, identified by its YouTube identifier. Two mechanisms keep repeated collection free of duplicates. In the database, the YouTube identifiers and the run links are unique; in the code, resources are written with \texttt{update\_\allowbreak or\_\allowbreak create}, which updates an existing row instead of inserting a new one. A test collects the same query twice, with different data in the second run, and verifies that each resource is stored once, that its values and retention timestamp were refreshed and that each run has its own links.

\subsection{Retention and purge}

The 30-day limit of the YouTube API Services Developer Policies \autocite{youtube_api_services_developer_policies} is enforced by a management command, \texttt{purge\_\allowbreak youtube\_\allowbreak data}. It deletes channels, videos and comments whose last retrieval is older than the retention period, and statistics snapshots by their capture time. The period is the setting \texttt{YOUTUBE\_\allowbreak RETENTION\_\allowbreak DAYS}: 30 days by default, as the policies require, and 0 to turn the purge off. Because the policies allow at most 30 days, a value above 30 or 0 produces a warning in Django's system checks at every start; the setting exists for a future change of the policy, not to keep API data longer. The deletion cascades to run links, replies, statistics snapshots and sentiment results, and the normalized text disappears with its comment because it is a column of the same row. Projects, saved queries, collection runs with their counters and parameters, request logs and processing reports are kept, because they describe the procedure and contain no YouTube data. When the purge removes the last resource linked to a run, it records the moment in the run's \texttt{data\_\allowbreak purged\_\allowbreak at} field, which the interface uses to explain why a run's data is no longer visible.

The whole purge runs in one transaction, so a failure leaves no partial deletion, and a \texttt{-{}-dry-run} option executes the same code and rolls the transaction back, so its reported counts are exact. The command runs every hour without an operator: the Docker Compose service \texttt{purge}, built from the same image as the API, repeats it in a shell loop, and the local start script does the same. It skips, with one line in its log, while any collection run is \texttt{running}, because a video refreshed by a running collection could otherwise still be deleted. The check and the purge are not one atomic step, so a run that starts between them can still overlap the purge; the outcome is at most a refreshed video deleted early, never corrupted data. A test takes a full copy of every table before and after a purge and checks that only the expected rows disappeared and that every remaining row is unchanged.
